\section{Rappels de 5ème}


\subsection{Vocabulaire}

\begin{exemplebox}
\includegraphics[
    width=\linewidth,
    keepaspectratio
]{images/droite_graduée.png}%
\end{exemplebox}
\begin{definitionbox}
\begin{itemize}
    \item L'ensemble des nombres \hlb{relatifs} est la réunion de l'ensemble
        des nombres \answer{positifs} et de l'ensemble des nombres
        \answer{négatifs}.
    \item Sur une droite graduée, tout point est repéré par un nombre relatif
        appelé \answer{abscisse}. 
        \begin{itemize}
            \item Sur l'exemple 1 : A a pour abscisse 3, on le note A(3) 
        \end{itemize}
    \item Pour un nombre donné, sa \answer{distance à zéro} est la distance
        entre le point d'abscisse 0 et le point d'abscisse le nombre donné.
        \begin{itemize}
            \item Pour (-4) et (+4) 1 : leur distance à 0 est 4 
        \end{itemize}
    \item Deux \answer{nombres opposés} ont la même distance à zéro et sont de
        signes contraire.
\end{itemize}
\end{definitionbox}

\begin{propbox}
    \bb{Règle 1} : Pour \bb{ajouter} deux termes de \hl{même signe}:
    \begin{itemize}
        \item On ajoute les distances à zéro
        \item On garde le même signe
    \end{itemize}
\end{propbox}


\begin{exemplebox}
    \begin{itemize}
        \item $A = (+5) + (+7) = \answer{12}$ 
        \item $B = (-5) + (-10) = \answer{15}$ 
    \end{itemize}
\end{exemplebox}

\begin{propbox}
    \bb{Règle 2} : Pour \bb{ajouter} deux termes de \hl{signes contraires}: on
    \bb{garde le signe du plus grand} et on soustrait leur \bb{différence}
\end{propbox}


\begin{exemplebox}
    \begin{itemize}
        \item $A = (5) + (-7) = \answer{-2}$ 
        \item $B = (-5) + (+10) = \answer{+5}$ 
    \end{itemize}
\end{exemplebox}


